\input{../tex-inputs/latex-leseansicht-vorspann}

\section[Arthur Schnitzler an Gustav Schwarzkopf, {[}zwischen 7. und 9. 11. 1902?{]}]{L04192 Arthur Schnitzler an Gustav Schwarzkopf, {[}zwischen 7. und 9. 11. 1902?{]}}
\nopagebreak\mylabel{L04192v}
\rehead{ }\normalsize\beginnumbering\briefempfaengerindex{Schwarzkopf, Gustav@\textsc{Schwarzkopf, Gustav}!zzzSchnitzler, Arthur@\emph{von Arthur Schnitzler}!1902-11-071@{{[}zwischen 7. und 9. 11. 1902?{]}}|(be}
\toendnotes[C]{\smallbreak\pagebreak[2]}
\correspDesc{Versand  durch Arthur Schnitzler am [zwischen 7. und 9. 11. 1902?] in Wien
\newline{}Erhalt  durch Gustav Schwarzkopf am [zwischen 7. und 9. 11. 1902?] in Wien}\toendnotes[C]{\smallbreak}
\Standort{CUL, Schnitzler, B 96.}
\physDesc{Karte, 91 Zeichen
\newline{}Handschrift: Bleistift, deutsche Kurrent}\toendnotes[C]{\smallbreak}
\pstart
           \noindent{}{\pb}lieber Guſtav, bitte \label{K_L04192-1v}\edtext{ko{\geminationm}en Sie Dinſtag}{\lemma{\textnormal{\emph{kommen Sie Dinstag}}}\Cendnote{\textnormal{Die Karte ist undatiert und 
                  nicht verlässlich zu datieren. Die Verwendung der Anrede (mit kleinem Anfangsbuchstaben bei »lieber«) und die Unterschrift spannen
                  einen Zeitraum von 1899 bis 1906 auf, in dem die Karte vielleicht gelaufen ist. Sofern Schwarzkopf\pwindex{Schwarzkopf, Gustav 7.\,11.\,1853 Wien – 13.\,11.\,1939 ebd.@\textsc{Schwarzkopf, Gustav} (7.\,11.\,1853 Wien – 13.\,11.\,1939 ebd.)|pwk}
               der Einladung nachkam und die Bedingung ›en famille‹ gleichblieb, lassen sich mit dem \emph{Tagebuch}\pwindex{Schnitzler, Arthur 15. 5. 1862 Wien – 21. 10. 1931 ebd.@\textsc{Schnitzler, Arthur} (15. 5. 1862 Wien – 21. 10. 1931 ebd.)!Tagebuch@\strich\emph{Tagebuch}|pwk} drei mögliche Termine
                  bestimmen: A. S.: \emph{Tagebuch}, 11. 11. 1902, A. S.: \emph{Tagebuch}, 26. 5. 1903 und A. S.: \emph{Tagebuch}, 15. 8. 1905. Eine Unterscheidung zwischen
               den Treffen ist nicht möglich, weswegen es hier nur unter starken Zweifeln beim ersten möglichen Termin eingeordnet wird.}}}\label{K_L04192-1}
      zu uns nachtmahlen; nur Familie.\pend
           
\pstart
           Herzlichſt Ihr{\\[\baselineskip]}\spacefill\mbox{Arthur Sch}\pend
           \leftskip=0em{}\selectlanguage{ngerman}\endnumbering\briefempfaengerindex{Schwarzkopf, Gustav@\textsc{Schwarzkopf, Gustav}!zzzSchnitzler, Arthur@\emph{von Arthur Schnitzler}!1902-11-071@{{[}zwischen 7. und 9. 11. 1902?{]}}|)be}\mylabel{L04192h}
\begin{anhang}
\end{anhang}\newcommand{\dateiname}{L04192}\newcommand{\titel}{Arthur Schnitzler an Gustav Schwarzkopf, [zwischen 7. und 9. 11. 1902?]}\newcommand{\editorInnen}{Herausgegeben von Jahnke, SelmaMüller, Martin Anton}\input{../tex-inputs/latex-leseansicht-abspann}
